\section{Opinión del analista}\label{sec:opinion}

El modelo tiene tres fortalezas que valen su costo computacional. Primero, revalorizar cada instrumento por
sus flujos completos, en vez de aproximar con duración o DV01, es lo que permite detectar que el riesgo de
tasa corta de la cartera óptima cambia de signo según la convención de proyección de cupones flotantes
(Sección~\ref{sec:sensibilidad}); una aproximación lineal habría escondido ese riesgo de modelo. Segundo, linealizar
el CVaR con la formulación de Rockafellar--Uryasev mantiene el problema como programa lineal, resoluble en
milisegundos, sin perder la coherencia entre el CVaR que entra en el objetivo y el que se reporta \emph{ex
post}. Tercero, la batería de nueve sensibilidades (tasas, precios sombra, costos, confianza, peso de la
cola, semilla, deriva, sesgo del estrés y regla de cupones) obliga a distinguir qué parte del resultado viene
de la optimización y qué parte viene de supuestos que el enunciado del caso no fija por completo.

Esa misma batería expone las debilidades del modelo. El conjunto de estrés castiga solo una dirección del
tipo de cambio: los cuatro escenarios deprecian el sol. Buena parte de la preferencia de la cartera
óptima por activos en dólares depende de esa asimetría (apartado~\ref{subsec:sesgo-estres}), y un comité de riesgos debería incorporar un
escenario de apreciación del sol al conjunto base, no solo como sensibilidad. El remuestreo por bloques
reproduce bien la razón de varianzas a 12 meses observada en la historia, pero esa historia cubre cinco años y
un solo ciclo de alza de tasas; la prueba de semillas alternativas muestra una dispersión de CVaR de
\PEN5.6 mm a \PEN7.7 mm que, sin ser grande, tampoco es despreciable frente al \PEN7.82 mm de la base. Y la
convención de proyección de cupones flotantes (nivel plano del factor de referencia frente a forward del
período) no está fijada de manera inequívoca por los datos del caso y cambia materialmente el resultado en
el escenario de estrés más severo; este informe la trata como una sensibilidad porque no hay evidencia
suficiente para preferir una convención sobre la otra, pero es la fuente de riesgo de modelo más grande de
todo el análisis.

Con más tiempo o más datos, extendería la historia de mercado más allá de cinco años para cubrir más de un
ciclo de tasas, y adoptaría como cartera base la variante que exige los límites de vencimiento también en el
horizonte: la solución base los incumple con certeza en el préstamo corto L01, y corregirlo cuesta solo
\PEN0.29 mm de resultado esperado neto. También agregaría de forma permanente al conjunto de estrés un
escenario de apreciación del sol, en vez de dejarlo como sensibilidad aparte, para que la cartera recomendada
no dependa de la dirección cambiaria de los cuatro escenarios que trae el caso.
